\documentclass[border=4pt]{standalone}
\usepackage{amsmath,amssymb,mathtools,xcolor,varwidth}
\begin{document}
\begin{varwidth}{28em}
\large\bfseries
Now we lay down Prop. VI's apparatus: $QT$ perpendicular to $SP$ and $LR$ parallel to $SP$,      meeting the tangent at $R$ and $Z$ and the circle at $L$.      The similar triangles $ZQR$, $ZTP$, $ZPA$ give $RP^2$ (that is $QR\cdot L$) to $QT^2$ as $AV^2$ to $PV^2$,      so multiplying by $SP^2/QR$ and letting $Q$ coincide with $P$ collapses the construction into the clean chord-cube.      Substituting into Prop. VI yields $F \propto \dfrac{1}{SP^2\cdot PV^3}$; and in the special case where $S$ sits at $O$,      the chord $PV$ becomes the diameter and the law reduces to the pure inverse-square $1/SP^2$. \textit{Q.E.D.}
\end{varwidth}
\end{document}
